\documentclass[12pt,a4paper]{article}
\usepackage{fontspec}
\setmainfont{cmunrm.otf}[Path=docs/fonts/,BoldFont=cmunbx.otf,ItalicFont=cmunti.otf,BoldItalicFont=cmunbi.otf]
\setmonofont{cmuntt.otf}[Path=docs/fonts/,BoldFont=cmuntb.otf,ItalicFont=cmunit.otf,BoldItalicFont=cmuntx.otf]
\usepackage[russian]{babel}
\usepackage{amsmath,amssymb,array,booktabs,caption,enumitem,fancyhdr,float,geometry,graphicx,hyperref,indentfirst,listings,longtable,tikz,xcolor}
\usetikzlibrary{arrows.meta,positioning}
\geometry{left=30mm,right=15mm,top=20mm,bottom=20mm}
\setlength{\parindent}{1.25cm}
\linespread{1.2}
\pagestyle{fancy}\fancyhf{}\fancyfoot[C]{\thepage}
\renewcommand{\headrulewidth}{0pt}\hypersetup{hidelinks}
\captionsetup{justification=centering}\setcounter{tocdepth}{2}
\definecolor{codebg}{RGB}{248,248,248}
\definecolor{codeframe}{RGB}{215,215,215}
\lstset{basicstyle=\ttfamily\small,backgroundcolor=\color{codebg},frame=single,
rulecolor=\color{codeframe},breaklines=true,showstringspaces=false,columns=fullflexible,keepspaces=true}
\lstset{aboveskip=5pt,belowskip=7pt}
\newcolumntype{P}[1]{>{\raggedright\arraybackslash}p{#1}}
\setlength{\emergencystretch}{3em}
\newcommand{\code}[1]{\texttt{\detokenize{#1}}}
\newcommand{\makelabtitle}[2]{
\pagenumbering{gobble}\hypersetup{pageanchor=false}
\begin{titlepage}\thispagestyle{empty}\begin{center}
{\small Федеральное государственное автономное образовательное учреждение высшего образования\\
«Национальный исследовательский университет ИТМО»}\\[0.8em]
{\small Факультет программной инженерии и компьютерной техники}
\vfill
{\Large\textbf{#1}}\\[0.6em]
{\large #2}\\[0.6em]
{\large по курсу «Системы искусственного интеллекта»}\\[0.6em]
{\large Модуль 1. Базы знаний и онтологии}
\vfill
\begin{flushright}\begin{tabular}{l}
Выполнили:\\Якименко Владислав, P3313\\Пахомов Игнат, P3313\\[0.8em]
Проверил:\\Жданов Андрей Дмитриевич
\end{tabular}\end{flushright}
\vfill Санкт-Петербург --- 2026
\end{center}\end{titlepage}\clearpage
\hypersetup{pageanchor=true}\tableofcontents\thispagestyle{empty}\clearpage\pagenumbering{arabic}}

\begin{document}
\makelabtitle{Лабораторная работа №1}{База знаний Prolog и онтология видеоигр}
\section{Цель и требования}
Создать базу знаний об играх и перевести её в онтологическую форму.
Требования: не менее 20 унарных фактов, 10--15 бинарных фактов, 5--7 правил,
запросы с переменными и логическими операциями. Для онтологии нужны
классы, свойства, индивиды, иерархия, ограничения и проверка вывода.
Работа включает обе части лабораторной 1 по выданному заданию.
\section{ЛР 1. База знаний Prolog}
\subsection{Предметная область и словарь}
Объекты представлены атомами в snake\_case: например, \code{portal_2},
\code{valve}. Переменные начинаются с заглавной буквы. Факт утверждает отношение,
а правило задаёт условия его доказательства. Запятая означает конъюнкцию,
точка с запятой --- дизъюнкцию. Отрицание \code{\+ Goal} успешно, когда
интерпретатор не смог доказать цель Goal.

\begin{table}[H]\centering
\caption{Факты исходной базы}\begin{tabular}{lr}\toprule
Предикат & Число фактов\\\midrule
\code{game/1} & 12\\
\code{single_player/1} & 11\\
\code{multiplayer/1} & 6\\
\code{story_driven/1} & 7\\
\code{difficult/1} & 5\\
\code{indie/1} & 4\\
\code{relaxing/1} & 3\\
\code{competitive/1} & 2\\
\code{free_to_play/1} & 1\\\midrule
Всего унарных & 51\\
\code{developed_by/2} & 10\\
\code{has_genre/2} & 5\\\bottomrule
\end{tabular}\end{table}

\subsection{Правила и порядок вычисления}
В файле games.pl факты сгруппированы по признакам и отношениям,
далее записаны правила. Группы фактов и основные правила снабжены комментариями.
Примеры запросов вынесены в queries.md и run\_examples.pl,
автоматические проверки --- в test\_games.pl; загрузка БЗ сама по себе
не запускает демонстрации или тесты.

Семь правил описывают сюжетные одиночные игры, дружелюбные многопользовательские,
напряжённые, сюжетные независимые, связь игры со студией, спокойные одиночные
и подходящие новичку игры. Пример композиции условий:
\begin{lstlisting}
recommended_for_beginner(Game) :-
    game(Game), single_player(Game),
    \+ difficult(Game), \+ competitive(Game).
\end{lstlisting}
Сначала выбирается конкретная игра, затем проверяются отрицания. Такая последовательность
существенна: проверка отрицания со свободной переменной спрашивала бы об отсутствии
любой сложной игры, а не о свойствах кандидата.
В \code{intense_game/1} игра Civilization VI удовлетворяет двум ветвям ИЛИ;
\code{setof/3} устраняет повторения в итоговом наборе.

\subsection{Отсечение и полнота}
Предикат \code{!} отсекает альтернативы, возникшие до него в пределах текущего вызова.
В запросе \code{story_driven(G), !, \+ developed_by(G,cd_projekt_red)}
первым выбирается The Witcher 3. Отрицание не выполняется, а возвращаться к
следующим сюжетным играм уже нельзя, поэтому ответ --- false.
Без отсечения находятся Portal 2, Hades, Hollow Knight, Celeste и Baldur's Gate 3.
Здесь cut меняет множество решений и не подходит для получения всех рекомендаций.

\subsection{Фактические ответы интерпретатора}
Ниже приведены десять запросов, выполненных скриптом сбора результатов.
\code{findall/3} собирает ответы в порядке поиска; \code{setof/3} возвращает
отсортированный набор. Списки совпадают с отдельными интерактивными решениями.
\par\medskip
\noindent\begin{minipage}{\linewidth}
\textbf{Запрос 1}
\begin{lstlisting}
?- game(minecraft).
true
\end{lstlisting}
\end{minipage}\par\smallskip

\noindent\begin{minipage}{\linewidth}
\textbf{Запрос 2}
\begin{lstlisting}
?- developed_by(portal_2,valve).
true
\end{lstlisting}
\end{minipage}\par\smallskip

\noindent\begin{minipage}{\linewidth}
\textbf{Запрос 3}
\begin{lstlisting}
?- findall(G,developed_by(G,valve),L).
[portal_2, counter_strike_2]
\end{lstlisting}
\end{minipage}\par\smallskip

\noindent\begin{minipage}{\linewidth}
\textbf{Запрос 4}
\begin{lstlisting}
?- findall(G,(single_player(G),story_driven(G)),L).
[the_witcher_3, cyberpunk_2077, portal_2, hades, hollow_knight,
celeste, baldurs_gate_3]
\end{lstlisting}
\end{minipage}\par\smallskip

\noindent\begin{minipage}{\linewidth}
\textbf{Запрос 5}
\begin{lstlisting}
?- findall(G,(game(G),(relaxing(G);free_to_play(G)),\+ competitive(G)),L).
[minecraft, stardew_valley]
\end{lstlisting}
\end{minipage}\par\smallskip

\noindent\begin{minipage}{\linewidth}
\textbf{Запрос 6}
\begin{lstlisting}
?- findall(G,friendly_multiplayer(G),L).
[portal_2, minecraft, stardew_valley, baldurs_gate_3]
\end{lstlisting}
\end{minipage}\par\smallskip

\noindent\begin{minipage}{\linewidth}
\textbf{Запрос 7}
\begin{lstlisting}
?- setof(G,intense_game(G),L).
[celeste, civilization_vi, counter_strike_2, doom_eternal, hades,
hollow_knight]
\end{lstlisting}
\end{minipage}\par\smallskip

\noindent\begin{minipage}{\linewidth}
\textbf{Запрос 8}
\begin{lstlisting}
?- findall(G,indie_gem(G),L).
[hades, hollow_knight, celeste]
\end{lstlisting}
\end{minipage}\par\smallskip

\noindent\begin{minipage}{\linewidth}
\textbf{Запрос 9}
\begin{lstlisting}
?- findall(G,recommended_for_beginner(G),L).
[the_witcher_3, cyberpunk_2077, portal_2, minecraft, stardew_valley,
baldurs_gate_3]
\end{lstlisting}
\end{minipage}\par\smallskip

\noindent\begin{minipage}{\linewidth}
\textbf{Запрос 10}
\begin{lstlisting}
?- findall(G,relaxing_solo_game(G),L).
[minecraft, stardew_valley, civilization_vi]
\end{lstlisting}
\end{minipage}\par\smallskip


\clearpage\section{Онтология в Protégé}
\subsection{Перевод фактов и правил}
Онтология содержит 21 именованный класс, 24 именованных индивида,
три объектных свойства, одно свойство данных, семь эквивалентностей и одно правило SWRL.
Из 24 индивидов 12 являются играми, 8 --- студиями и 4 --- жанрами.
Пространство имён: \url{http://example.org/itmo/games\#}.
OWL и Turtle представляют один RDF-граф из 329 троек; их изоморфность проверяется скриптом.

\begin{table}[H]\centering\small
\caption{Трассируемость Prolog и OWL}
\begin{tabular}{P{6.3cm}P{8.0cm}}\toprule
Prolog & OWL\\\midrule
\code{game(G)} & принадлежность VideoGame\\
\code{single_player(G)}, \code{multiplayer(G)} & SinglePlayerGame, MultiplayerGame\\
\code{story_driven(G)}, \code{difficult(G)} & StoryDrivenGame, DifficultGame\\
\code{indie(G)}, \code{relaxing(G)} & IndieGame, RelaxingGame\\
\code{competitive(G)}, \code{free_to_play(G)} & CompetitiveGame, FreeToPlayGame\\
\code{developed_by(G,S)} & developedBy(G,S), обратное developedGame(S,G)\\
\code{has_genre(G,T)} & hasGenre(G,T)\\
\code{solo_story_game(G)} & SoloStoryGame: SinglePlayerGame and StoryDrivenGame\\
\code{friendly_multiplayer(G)} & MultiplayerGame and NonCompetitiveGame\\
\code{intense_game(G)} & DifficultGame or CompetitiveGame\\
\code{indie_gem(G)} & IndieGame and StoryDrivenGame\\
\code{relaxing_solo_game(G)} & RelaxingGame and SinglePlayerGame\\
\code{recommended_for_beginner(G)} & SinglePlayerGame and NonDifficultGame and NonCompetitiveGame\\
\code{studio_game(G,S)} & StudioGame --- проекция на G; S извлекается через developedBy\\\bottomrule
\end{tabular}\end{table}

Двухаргументный предикат нельзя полностью заменить одним классом: StudioGame
показывает наличие разработчика, но для получения пары игра--студия необходимо
объектное свойство. В исходном наборе десять явно известных таких пар.
По семи остальным унарным результатам, включая отдельную проверку SWRL,
множества экземпляров сопоставляются с результатами Prolog автоматически.

\subsection{Иерархия и свойства}
\begin{figure}[H]\centering
\begin{tikzpicture}[every node/.style={draw,rounded corners,font=\small,align=center},>=Stealth]
\node (game) at (0,0) {VideoGame};
\node (single) at (-2,-1.3) {SinglePlayerGame};
\node (story) at (2,-1.3) {StoryDrivenGame};
\node (solo) at (0,-2.6) {SoloStoryGame};
\node (studio) at (5,0) {Studio};
\draw[->] (single)--node[draw=none,left]{\scriptsize is-a}(game);
\draw[->] (story)--(game);
\draw[->] (solo)--(single);\draw[->] (solo)--(story);
\draw[->] (game)--node[draw=none,above]{\scriptsize developedBy}(studio);
\end{tikzpicture}
\caption{Фрагмент иерархии и связи между играми и студиями}
\end{figure}
Классы VideoGame, Studio и Genre попарно несовместимы. При этом SinglePlayerGame
и MultiplayerGame совместимы: Portal 2 относится к обоим. Domain свойства developedBy
равен VideoGame, range --- Studio; у developedGame направления обратные.
У hasGenre domain VideoGame, range Genre. Domain и range служат аксиомами вывода
типов, а не проверками ввода в стиле реляционной базы.

Добавлен DataProperty title с диапазоном xsd:string. У VideoGame задано
title exactly 1 xsd:string, у StudioGame --- developedBy min 1 Studio.
Названия являются дополнительными данными отображения относительно исходного Prolog.
Минимальная кардинальность в OWL допускает неизвестный объект; отсутствие
явно записанного значения само по себе не делает онтологию противоречивой.
В отличие от этого, два разных строковых title нарушают exactly 1.

\subsection{Открытый мир и SWRL}
Prolog трактует отсутствие доказательства как неудачу цели. OWL не выводит
отрицание из отсутствия факта. Поэтому для данного конечного набора явно
заданы NonCompetitiveGame и NonDifficultGame, несовместимые с положительными
категориями. Эти утверждения добавлены при переводе и требуют пересмотра
при изменении набора фактов.

Для отдельной демонстрации SWRL задано правило:
\begin{lstlisting}
SinglePlayerGame(?g) ^ StoryDrivenGame(?g)
  -> RuleVerifiedStoryGame(?g)
\end{lstlisting}
Этот класс не имеет Equivalent To, поэтому его семь экземпляров демонстрируют
именно срабатывание правила. HermiT проверил правило для именованных индивидов.
Остальные логические определения выражены штатными конструкциями OWL.

\subsection{Компетентностные запросы и проверка}
В Protégé файл открывается через File → Open. После запуска Reasoner → HermiT →
Start reasoner на вкладке DL Query выбирается Instances и вводятся имена ниже.
При отображении prefixed names используется ведущий символ двоеточия.

\begin{enumerate}[leftmargin=*]
\item \code{SoloStoryGame}: baldurs\_gate\_3, celeste, cyberpunk\_2077, hades,
hollow\_knight, portal\_2, the\_witcher\_3.
\item \code{FriendlyMultiplayerGame}: baldurs\_gate\_3, minecraft, portal\_2, stardew\_valley.
\item \code{IntenseGame}: celeste, civilization\_vi, counter\_strike\_2,
doom\_eternal, hades, hollow\_knight.
\item \code{IndieGem}: celeste, hades, hollow\_knight.
\item \code{BeginnerRecommendedGame}: baldurs\_gate\_3, cyberpunk\_2077,
minecraft, portal\_2, stardew\_valley, the\_witcher\_3.
\item \code{RuleVerifiedStoryGame}: тот же набор из семи игр, что у SoloStoryGame.
\end{enumerate}

Наборы выше получены HermiT в автоматическом запуске, а не переписаны
из предполагаемого результата. Файл \code{verification.json} фиксирует вывод.
Для визуального просмотра можно включить OntoGraf и раскрыть VideoGame;
рисунок в отчёте показывает фрагмент тех же отношений.
Дополнительно SPARQL с FILTER NOT EXISTS выбирает сюжетные игры без записанной
связи с Valve: baldurs\_gate\_3, celeste, cyberpunk\_2077, hades, hollow\_knight,
the\_witcher\_3. Это проверка отсутствия тройки в графе, а не OWL-доказательство
отрицания. Запрос DL \code{not (developedBy value valve)} в общем случае
не эквивалентен такому фильтру.

\section{Тестирование и вывод}
Все 13 тестов Prolog прошли в SWI-Prolog 10.0.2. Выполнены десять запросов,
приведённых в листингах. HermiT подтвердил консистентность онтологии,
семь наборов выведенных типов и обратное свойство. Негативные проверки
disjoint, типа данных и кардинальности обнаружили противоречия.

Исходная БЗ удовлетворяла требованиям по объёму, но онтология нуждалась
в добавлении DataProperty, кардинальностей и SWRL. Эти дополнения выполнены.
Свойства игр приняты как учебные факты; жанровые данные неполны.
Отрицательные классы OWL заданы явно для согласования с конечной Prolog-базой.
При расширении базы их необходимо актуализировать.

В результате получена воспроизводимая модель предметной области, пригодная
для последующего использования в рекомендательной системе. Команды запуска,
запросы и протокол проверки включены в папки lab\_1 и reports.

\section{Источники}
Методическое пособие «Системы искусственного интеллекта», ИТМО, 2024,
с. 56--57, приложение 2: \url{https://books.ifmo.ru/file/pdf/3308.pdf}.
Задание «Системы ИИ.docx», лабораторная 1, части 1 и 2.
\end{document}
